\section{Conclusion}

This submission version answers two core questions, each with a qualification established by the extended falsification checks in Section~6. Within-country increases in the composite AI/digital index are consistently associated with higher TFP, but decomposition shows this association is carried by the digital-infrastructure component rather than the innovation/patent component, and a reverse-causality check finds that lagged TFP itself positively predicts current AI adoption. The association with contemporaneous GDP growth is positive but smaller, and the same reverse-causality caveat applies to its timing interpretation.

These findings are stable across principal robustness and sensitivity checks, including two-way fixed effects, lag structures, alternative clustering, and Driscoll-Kraay inference. The direction and broad magnitude of the digital/AI-TFP association are therefore unlikely to be artifacts of one narrow specification choice, but its construct-validity and timing limits are now characterized directly rather than only discussed in principle.

The contribution is empirical and methodological: stable findings across a structured robustness framework, produced through a reproducible pipeline with transparent reporting. The interpretive boundary remains explicit: these are robust conditional associations between digital/AI capacity and productivity, not definitive or AI-specific causal effects.

For policy, the results are consistent with a complementary-capabilities perspective. AI adoption appears more likely to translate into measurable productivity gains where countries also sustain investments in human capital, digital infrastructure, and implementation capacity. For research, the priority is stronger causal identification and richer analysis of heterogeneous effects.

Compared with related evidence, this pattern supports complementarity-based diffusion narratives but is more conservative about near-term aggregate growth effects, emphasizing gradual macro transmission rather than immediate transformation.

An online appendix provides full regression outputs, extended diagnostics, and additional figures.
